% !PRINTED 222
\weberpage{222}

%<<CONT
whose roots are the $x_{h,0}$. The influence of a substitution of the group
$\mathfrak{A}_{1,0}$
\[
\begin{pmatrix} a, & b \\ 0, & \partial \end{pmatrix}
\]
upon $x_{h,0}$ consists in this, that $x_{h,0}$ passes over into $x_{\partial h,0}$, and the
group of this equation of degree $\varphi(n)$ therefore consists of the permutations
\[
(x_{h,0},\ x_{\partial h,0}),
\]
in which $\partial$ may be any number whatever coprime to $n$. This group is
an \gesperrt{Abelian} one, and therefore the roots $x_{h,0}$ are, \gesperrt{after the adjunction of $\xi$,
to be determined algebraically by radicals.}

If, however, we adjoin not merely one, but all the roots $\xi$ of the
transformation equation, or, what comes to the same thing, the three belonging to
$R_{1,0}$, $R_{0,1}$, $R_{1,1}$, then the group of the
division equation reduces to $\mathfrak{A}_0$, and if we adjoin in addition $n$th roots of unity,
to $\mathfrak{B}_0$. The group $\mathfrak{B}_0$ is an \gesperrt{Abelian} one of
order $2^k$ which contains only elements of order $= 2$. In consequence the
\gesperrt{division equation is solvable by square roots.}

In order to find the form of this solution, let one put
\[
n = p^\pi p'^{\pi'} \ldots,
\]
in which $p, p', \ldots$ are primes different from one another, and $\pi, \pi', \ldots$
are positive exponents. Let the numbers $c, c', \ldots$ be determined from the congruences
\[
\begin{aligned}
c &\equiv -1\ (\mathrm{mod}\ p^\pi), & c' &\equiv -1\ (\mathrm{mod}\ p'^{\pi'}) \ldots \\
c &\equiv +1\ (\mathrm{mod}\ np^{-\pi}), & c' &\equiv +1\ (\mathrm{mod}\ np'^{-\pi'}) \ldots
\end{aligned}
\]
then one obtains every solution $\alpha$ of the congruence
\[
\alpha^2 \equiv 1\ (\mathrm{mod}\ n),
\]
and each only once, in the form
\weq{12}{\alpha \equiv c^\varepsilon c'^{\varepsilon'} \ldots \ (\mathrm{mod}\ n),}
when $\varepsilon, \varepsilon', \ldots$ take the values 0, 1.

Now if $\Omega$ is any one of the values
\[
\frac{4\mu K + 4\mu' i K'}{n},
\]
%CONT>>
